%% physical_feasibility.tex — PCSS Theory: L3 Physical Feasibility
%% Paper F: "Physically-Constrained Semantic Security Theory for IoT Mesh Routing"
%% Issue #81 (F3)
%%
%% Depends on: definitions.tex (Definitions 1–6), theorem1_proof.tex (Theorem 1)
%% Referenced in: SP-PaperF.tex §3.3
%%
%% Required preamble packages:
%%   \usepackage{amsthm, amsmath, amssymb, mathtools}
%%
%% Theorem environments (declare in preamble if not already present):
%%   \theoremstyle{definition}
%%   \newtheorem{definition}{Definition}[section]
%%   \newtheorem{remark}[definition]{Remark}
%%   \newtheorem{example}[definition]{Example}
%%   \theoremstyle{plain}
%%   \newtheorem{theorem}{Theorem}[section]

% -----------------------------------------------------------------------
\subsection{L3 — Physical Feasibility}
\label{subsec:l3-feasibility}
% -----------------------------------------------------------------------

Theorem~\ref{thm:sbc-sufficiency} establishes that $\mathrm{SBC}(P)$
is \emph{sufficient} to exclude a Dolev-Yao attacker from the routing
table, provided authentication can be applied to every message in
$\mathrm{DC}(P)$.  In the classical cryptographic model, computational
resources are unlimited: a protocol is simply labelled ``correct'' or
``incorrect'' without regard to whether it can be executed at all.

IoT radio nodes operate under two hard physical ceilings that the
standard model ignores:
\begin{enumerate}[label=(\roman*)]
  \item \textbf{Duty-cycle limit.}  ETSI EN~300~220 and FCC Part~15
        restrict the fraction of time a LoRa transmitter may occupy the
        channel to $\delta_{\max} = 1\%$ in the 868~MHz sub-band.  A node
        that authenticates and re-transmits control messages at rate $\rho$
        with airtime $t_{\mathrm{tx}}$ per message must satisfy
        $\rho \cdot t_{\mathrm{tx}} \le \delta_{\max}$; exceeding this is
        illegal and causes interference.

  \item \textbf{Battery energy budget.}  A node drawing current $I$ for
        time $t$ consumes $I \cdot t$ mAh from a finite battery of
        capacity $B_{\mathrm{mAh}}$.  Authentication adds a non-trivial
        computational current $I_{\mathrm{hmac}}$ for duration
        $t_{\mathrm{hmac}}$ per message; at rate $\rho$ these costs
        accumulate and determine the maximum protocol lifetime.
\end{enumerate}

We formalise both constraints via three definitions and then prove
Theorem~2, which gives a closed-form \emph{maximum safe authentication
rate} $\rho^*$ for any device--deployment pair.

% -----------------------------------------------------------------------
\begin{definition}[Environment Model]
\label{def:environment}
An \emph{environment model} is a tuple
\[
  E \;=\;
  \bigl(
    B_{\mathrm{mAh}},\;
    I_{\mathrm{tx}},\; t_{\mathrm{tx}},\;
    \delta_{\max},\;
    I_{\mathrm{hmac}},\; t_{\mathrm{hmac}},\;
    L_{\mathrm{days}}
  \bigr),
\]
where:
\begin{itemize}
  \item $B_{\mathrm{mAh}} \in \mathbb{R}_{>0}$ is the \emph{battery
        capacity} of the node (milliampere-hours);

  \item $I_{\mathrm{tx}} \in \mathbb{R}_{>0}$ is the \emph{transmit
        current} drawn by the radio during packet transmission (mA);

  \item $t_{\mathrm{tx}} \in \mathbb{R}_{>0}$ is the \emph{on-air time}
        per control-plane packet at the chosen spreading factor and
        bandwidth (ms);

  \item $\delta_{\max} \in (0, 1)$ is the \emph{maximum duty-cycle
        fraction} mandated by the applicable radio regulation
        (dimensionless; default 0.01 for ETSI 868~MHz);

  \item $I_{\mathrm{hmac}} \in \mathbb{R}_{>0}$ is the \emph{HMAC
        computation current} drawn by the MCU while executing
        HMAC-SHA256 (mA);

  \item $t_{\mathrm{hmac}} \in \mathbb{R}_{>0}$ is the \emph{HMAC
        computation time} per message on the target MCU (ms); and

  \item $L_{\mathrm{days}} \in \mathbb{R}_{>0}$ is the \emph{target
        node lifetime} (days).
\end{itemize}
All parameters are deterministic constants characterising a specific
device and deployment; stochastic channel effects are absorbed into the
choice of $t_{\mathrm{tx}}$ (computed for the worst-case payload length).
\end{definition}

% -----------------------------------------------------------------------
\begin{definition}[Authentication Cost Function]
\label{def:auth-cost}
Given an environment $E$, the \emph{authentication cost function}
$C_{\mathrm{auth}}(E)$ is the total energy expended per authenticated
control-plane message, expressed in milliampere-hours:
\[
  C_{\mathrm{auth}}(E)
  \;=\;
  \frac{I_{\mathrm{hmac}} \cdot t_{\mathrm{hmac}}
        \;+\; I_{\mathrm{tx}} \cdot t_{\mathrm{tx}}}{3\,600\,000}
  \quad \text{[mAh/msg]},
\]
where the factor $3\,600\,000$ converts milliampere-milliseconds to
milliampere-hours
($1\ \mathrm{mAh} = 3\,600\ \mathrm{s} \times 1000\ \mathrm{ms/s}
  = 3\,600\,000\ \mathrm{mA{\cdot}ms}$).

The first term $I_{\mathrm{hmac}} \cdot t_{\mathrm{hmac}}$ captures the
MCU energy for computing the HMAC tag; the second term
$I_{\mathrm{tx}} \cdot t_{\mathrm{tx}}$ captures the radio-transmit energy
for the authenticated packet.  Reception and verification energy is
conservatively neglected (receive current is typically $\le 15\%$ of
transmit current for LoRa modules).
\end{definition}

% -----------------------------------------------------------------------
\begin{definition}[Physical Feasibility Predicate]
\label{def:feasibility}
Let $P$ be a routing protocol, $E$ an environment model, and
$\rho_{\mathrm{DC}}(P) \ge 0$ the \emph{required authentication rate},
i.e., the minimum number of $\mathrm{DC}(P)$ messages per second that a
node must authenticate and transmit to participate in the protocol.
Define two rate bounds:
\begin{align}
  \rho_{\delta}(E) &\;=\; \frac{\delta_{\max}}{t_{\mathrm{tx}}/1000}
    \qquad \text{[pkt/s], duty-cycle bound,}
    \label{eq:rho-delta}\\[4pt]
  \rho_{B}(E)       &\;=\; \frac{B_{\mathrm{mAh}}}
                              {C_{\mathrm{auth}}(E) \cdot 86400 \cdot L_{\mathrm{days}}}
    \qquad \text{[pkt/s], battery-lifetime bound,}
    \label{eq:rho-B}
\end{align}
where $t_{\mathrm{tx}}/1000$ converts milliseconds to seconds in
\eqref{eq:rho-delta}, and $86400$ converts days to seconds in
\eqref{eq:rho-B}.
The \emph{maximum feasible authentication rate} is
\[
  \rho^*(E) \;=\; \min\!\bigl(\rho_{\delta}(E),\; \rho_{B}(E)\bigr).
\]
The \emph{physical feasibility predicate} $\Phi(P, E)$ is:
\[
  \Phi(P, E) \;=\; \mathbf{1}\!\left[\,\rho_{\mathrm{DC}}(P) \;\le\; \rho^*(E)\,\right].
\]
$\Phi(P, E) = 1$ asserts that the protocol can sustain full-rate
authentication of $\mathrm{DC}(P)$ messages within both the
duty-cycle budget and the battery-lifetime budget of $E$.
$\Phi(P, E) = 0$ asserts that at least one physical constraint is
violated.
\end{definition}

% -----------------------------------------------------------------------
\begin{theorem}[Physical Feasibility Bound]
\label{thm:feasibility-bound}
For every environment model $E$ there exists a unique $\rho^*(E) > 0$
such that
\[
  \Phi(P, E) = 1
  \quad \Longleftrightarrow \quad
  \rho_{\mathrm{DC}}(P) \;\le\; \rho^*(E),
\]
with the closed form
\[
  \rho^*(E) \;=\;
  \min\!\left(
    \frac{\delta_{\max}}{t_{\mathrm{tx}}/1000},\;
    \frac{B_{\mathrm{mAh}}}{C_{\mathrm{auth}}(E) \cdot 86400 \cdot L_{\mathrm{days}}}
  \right).
\]
\end{theorem}

\begin{proof}
By Definitions~\ref{def:environment}--\ref{def:feasibility}, all
parameters satisfy strict positivity ($B_{\mathrm{mAh}},\;
I_{\mathrm{tx}},\; t_{\mathrm{tx}},\; \delta_{\max},\;
I_{\mathrm{hmac}},\; t_{\mathrm{hmac}},\; L_{\mathrm{days}} > 0$).
Therefore:
\begin{enumerate}[label=(\roman*)]
  \item $\rho_{\delta}(E) = \delta_{\max}/(t_{\mathrm{tx}}/1000) > 0$
        because $\delta_{\max} > 0$ and $t_{\mathrm{tx}} > 0$.

  \item $C_{\mathrm{auth}}(E) > 0$ because
        $I_{\mathrm{tx}} \cdot t_{\mathrm{tx}} > 0$, hence
        $\rho_{B}(E) = B_{\mathrm{mAh}} /
        (C_{\mathrm{auth}}(E) \cdot 86400 \cdot L_{\mathrm{days}}) > 0$.

  \item $\rho^*(E) = \min(\rho_{\delta}, \rho_{B}) > 0$ as the minimum
        of two positive reals.
\end{enumerate}
Uniqueness is immediate since $\rho^*(E)$ is fully determined by $E$.
The biconditional follows directly from Definition~\ref{def:feasibility}:
$\Phi(P, E) = 1$ iff $\rho_{\mathrm{DC}}(P) \le \rho^*(E)$, and
$\Phi(P, E) = 0$ otherwise.
\hfill$\square$
\end{proof}

% -----------------------------------------------------------------------
\begin{remark}[Why $\Phi$ matters: degradation of SBC under infeasibility]
\label{rem:phi-degradation}
Theorem~\ref{thm:sbc-sufficiency} requires $\mathrm{SBC}(P)$ to hold, which
in turn requires every message in $\mathrm{DC}(P)$ to be authenticated at
the rate $\rho_{\mathrm{DC}}(P)$ at which it is produced.  If
$\Phi(P, E) = 0$, the protocol is \emph{cryptographically correct} (SBC
holds in the abstract model) but \emph{physically infeasible}: the node
cannot sustain full-rate HMAC computation and transmission without either
violating its duty-cycle allocation or exhausting its battery before
$L_{\mathrm{days}}$.

The practical consequence is that the implementation must reduce the
effective authentication rate $\rho' < \rho_{\mathrm{DC}}(P)$, for example
by authenticating only a fraction of $\mathrm{DC}(P)$ messages or by
increasing the inter-message interval.  This introduces a \emph{replay
window} of duration $1/\rho'$ seconds during which a stale authenticated
message can be re-injected without violating the freshness check
(condition~\ref{sbc:fresh} of Definition~\ref{def:sbc}).  The effective
security guarantee degrades from the unconditional exclusion of
Theorem~\ref{thm:sbc-sufficiency} to a probabilistic one parameterised
by the replay window length.

The classical Dolev-Yao model has no mechanism to express this
degradation: an attacker is either computationally bounded or not.
The PCSS framework makes the degradation explicit through $\Phi(P, E)$
and $\rho^*(E)$, enabling quantitative protocol design: choose the
control-message rate $\rho_{\mathrm{DC}}(P) \le \rho^*(E)$ so that
$\mathrm{SBC}(P)$ is simultaneously cryptographically sound and
physically sustainable.
\end{remark}

% -----------------------------------------------------------------------
\begin{example}[EoRa-PI instantiation of $\rho^*$]
\label{ex:eora-pi}
We instantiate Definition~\ref{def:feasibility} for the Ebyte EoRa-PI
ESP32-S3 LoRa node (the hardware platform used in Papers~B and~C) with
the following measured and regulation-fixed parameters:
\[
  E_{\mathrm{EoRa}} \;=\;
  \bigl(
    \underbrace{3000}_{\textstyle B_{\mathrm{mAh}}},\;
    \underbrace{120}_{\textstyle I_{\mathrm{tx}}},\;
    \underbrace{1692}_{\textstyle t_{\mathrm{tx}}},\;
    \underbrace{0.01}_{\textstyle \delta_{\max}},\;
    \underbrace{80}_{\textstyle I_{\mathrm{hmac}}},\;
    \underbrace{0.5}_{\textstyle t_{\mathrm{hmac}}},\;
    \underbrace{365}_{\textstyle L_{\mathrm{days}}}
  \bigr),
\]
where $t_{\mathrm{tx}} = 1692\ \mathrm{ms}$ is the measured airtime for a
142-byte LoRa frame at SF9/BW125 (the LoRaMesher Hello message size), and
$I_{\mathrm{hmac}} = 80\ \mathrm{mA}$, $t_{\mathrm{hmac}} = 0.5\ \mathrm{ms}$
are measured on the ESP32-S3 executing HMAC-SHA256 in software.

\medskip
\noindent\textbf{Step~1: Authentication cost per message (Definition~\ref{def:auth-cost}).}
\begin{align*}
  C_{\mathrm{auth}}(E_{\mathrm{EoRa}})
  &= \frac{80 \times 0.5 \;+\; 120 \times 1692}{3\,600\,000} \\
  &= \frac{40 \;+\; 203\,040}{3\,600\,000} \\
  &= \frac{203\,080}{3\,600\,000}
  \;\approx\; 5.641 \times 10^{-2}\ \mathrm{mAh/msg}.
\end{align*}

\medskip
\noindent\textbf{Step~2: Duty-cycle bound $\rho_\delta$.}
\begin{align*}
  \rho_{\delta}(E_{\mathrm{EoRa}})
  &= \frac{0.01}{1692/1000}
   = \frac{0.01}{1.692}
   \approx 5.91 \times 10^{-3}\ \mathrm{pkt/s}
   \;\approx 1\ \mathrm{pkt\ per\ 169\ s.}
\end{align*}

\medskip
\noindent\textbf{Step~3: Battery-lifetime bound $\rho_B$.}
\begin{align*}
  \rho_{B}(E_{\mathrm{EoRa}})
  &= \frac{3000}{5.641\times10^{-2} \times 86400 \times 365} \\
  &= \frac{3000}{1\,778\,981}
   \;\approx\; 1.69 \times 10^{-3}\ \mathrm{pkt/s}
   \;\approx 1\ \mathrm{pkt\ per\ 593\ s\ (\approx 9.9\ min).}
\end{align*}

\medskip
\noindent\textbf{Step~4: Maximum feasible rate.}
\[
  \rho^*(E_{\mathrm{EoRa}})
  \;=\; \min\!\bigl(5.91 \times 10^{-3},\; 1.69 \times 10^{-3}\bigr)
  \;=\; 1.69 \times 10^{-3}\ \mathrm{pkt/s}.
\]
The binding constraint is the \textbf{battery-lifetime bound}, not the
duty-cycle limit: at $\rho^*$ the node consumes approximately
$5.641 \times 10^{-2} \times 1.69 \times 10^{-3} \times 86400 \times 365
\approx 3000\ \mathrm{mAh}$ over 365 days, exactly exhausting the
battery.  The duty-cycle budget is only $28.6\%$ consumed
($1.69 \times 10^{-3} / 5.91 \times 10^{-3} \approx 0.286$), leaving
headroom for data traffic.

\medskip
\noindent\textbf{Consequence for LoRaMesher ($P_{\mathrm{LM}}$).}\;
The default LoRaMesher Hello interval is 30~s, corresponding to
$\rho_{\mathrm{DC}}(P_{\mathrm{LM}}) = 1/30 \approx 3.33 \times 10^{-2}\ \mathrm{pkt/s}$.
Since $3.33 \times 10^{-2} \gg 1.69 \times 10^{-3}$, we have
$\Phi(P_{\mathrm{LM}}, E_{\mathrm{EoRa}}) = 0$: full-rate HMAC
authentication of every Hello message is \emph{physically infeasible} on
this hardware at the default interval.  The implication (Remark~\ref{rem:phi-degradation})
is that the deployment must either (a)~increase the Hello interval to
$\ge 593\ \mathrm{s}$ (reducing routing agility), or (b)~authenticate a
fraction of messages (weakening SBC2 freshness), or (c)~use a higher-capacity
battery.  The PCSS framework makes this trade-off explicit and
quantitatively precise.
\end{example}

\begin{remark}[Scope of $E$: per-node vs.\ network-level]
\label{rem:scope}
Definition~\ref{def:environment} models a \emph{single node}.  In a
heterogeneous deployment, each $v_i$ may have its own $E_i$ and hence
its own $\rho^*(E_i)$.  The network-wide feasibility bound is
$\rho^*_{\mathrm{net}} = \min_{v_i \in V} \rho^*(E_i)$: the most
energy-constrained node determines the maximum protocol-wide
authentication rate.  This motivates per-node duty-cycle adaptation
mechanisms in multi-hop LoRa mesh protocols, which fall outside the scope
of the present paper.
\end{remark}
